\documentclass[11pt]{article}
\usepackage[margin=1in]{geometry}
\usepackage{array}
\usepackage{enumitem}
\usepackage{xcolor}
\usepackage{hyperref}
\usepackage{listings}

\hypersetup{colorlinks=true, linkcolor=blue, urlcolor=blue}
\setlength{\parindent}{0pt}
\setlength{\parskip}{6pt}
\setlist[itemize]{topsep=3pt,itemsep=2pt,leftmargin=*}
\setlist[enumerate]{topsep=3pt,itemsep=4pt,leftmargin=*}

\lstset{
  basicstyle=\ttfamily\small,
  breaklines=true,
  frame=single,
  columns=fullflexible,
  keepspaces=true,
  showstringspaces=false,
  backgroundcolor=\color{gray!6}
}

\definecolor{predictcolor}{RGB}{255,242,200}
\definecolor{discusscolor}{RGB}{214,234,255}
\definecolor{funfactcolor}{RGB}{222,247,222}
\definecolor{debatecolor}{RGB}{255,221,221}

\newcommand{\calloutbox}[3]{%
  \vspace{4pt}\noindent
  \colorbox{#1}{\parbox{0.96\linewidth}{\textbf{#2}}}\\[2pt]
  \noindent\fbox{\begin{minipage}{0.94\linewidth}\vspace{2pt}#3\vspace{2pt}\end{minipage}}
  \vspace{6pt}
}
\newcommand{\predictbox}[1]{\calloutbox{predictcolor}{PREDICTION TIME}{#1}}
\newcommand{\discussbox}[1]{\calloutbox{discusscolor}{HUDDLE UP}{#1}}
\newcommand{\funfact}[1]{\calloutbox{funfactcolor}{FUN FACT}{#1}}
\newcommand{\debatebox}[1]{\calloutbox{debatecolor}{DEBATE CORNER}{#1}}
\newcommand{\claimbox}[1]{\calloutbox{funfactcolor}{CLAIM, EVIDENCE, CONTEXT}{#1}}

\newcommand{\writebox}[1]{
  \vspace{3pt}
  \noindent\fbox{
    \begin{minipage}{0.96\linewidth}
      \textbf{#1}\\[12pt]
      \rule{\linewidth}{0.4pt}\\[14pt]
      \rule{\linewidth}{0.4pt}\\[14pt]
      \rule{\linewidth}{0.4pt}\\[14pt]
    \end{minipage}
  }
  \vspace{6pt}
}

\newcommand{\shortwritebox}[1]{
  \vspace{3pt}
  \noindent\fbox{
    \begin{minipage}{0.96\linewidth}
      \textbf{#1}\\[12pt]
      \rule{\linewidth}{0.4pt}\\[14pt]
    \end{minipage}
  }
  \vspace{6pt}
}

\begin{document}

\begin{center}
  {\LARGE MA 213 Week 9: Lab 6}\\[3pt]
  {\large Confidence Intervals and Hypothesis Tests}
\end{center}

\begin{enumerate}
  \item \textbf{Your name:} \underline{\hspace{0.3\textwidth}} \hfill \textbf{BU ID:} \underline{\hspace{0.25\textwidth}}
  \item \textbf{Partner's name:} \underline{\hspace{0.3\textwidth}} \hfill \textbf{BU ID:} \underline{\hspace{0.25\textwidth}}
\end{enumerate}

\textbf{Big idea:} Confidence intervals and hypothesis tests answer different but related questions. In this lab, you will simulate confidence interval coverage and manually build a chi-square test of independence.
\textbf{Do not use GenAI tools for this lab.} The point is to practice your own inference reasoning, coding skills, and interpretation.

\textbf{Files used:} No external data file is needed. You will generate one population and use R's built-in \texttt{Titanic} table.

\textbf{Skills practiced:} \texttt{sample()}, \texttt{qnorm()}, \texttt{pchisq()}, \texttt{table()}, \texttt{addmargins()}, \texttt{as.data.frame()}, \texttt{group\_by()}, and \texttt{summarize()}.

\textbf{Your mission:} Explain what 95\% confidence means through simulation, then test whether two categorical variables appear independent.

\section*{Icebreaker}

Before opening R: introduce yourself to your partner. Then answer: Which sounds more convincing to you, ``we estimate the proportion is between 60\% and 64\%'' or ``we reject the null hypothesis''? Why?

\shortwritebox{My answer and my partner's answer:}

\section*{Getting Started: One Proportion}

Suppose a survey asks Boston residents whether they own a cell phone. Out of 200 randomly selected residents, 127 say yes.

\begin{lstlisting}[language=R]
library(dplyr)
library(ggplot2)

x <- 127
n <- 200
p_hat <- x / n

p_hat
\end{lstlisting}

\section*{Question 1. Construct and interpret a confidence interval}

Construct a 95\% confidence interval for the population proportion of Boston residents who own cell phones.

\[
\hat{p} \pm z^* \sqrt{\frac{\hat{p}(1-\hat{p})}{n}}
\]

\begin{lstlisting}[language=R]
alpha <- 0.05
z_star <- qnorm(1 - alpha / 2)

SE <- sqrt(p_hat * (1 - p_hat) / n)
CI <- p_hat + c(-1, 1) * z_star * SE

z_star
SE
CI
\end{lstlisting}

\writebox{Interpret the confidence interval in context. Avoid saying there is a 95\% probability that this one interval contains the true proportion.}

\section*{Question 2. Test a hypothesis using the same sample}

Now test whether the proportion of Boston residents who own cell phones is different from 50\%.

\[
H_0: p = 0.50
\qquad
H_a: p \neq 0.50
\]

\begin{lstlisting}[language=R]
p0 <- 0.50

SE_null <- sqrt(p0 * (1 - p0) / n)
z_stat <- (p_hat - p0) / SE_null
p_value <- 2 * (1 - pnorm(abs(z_stat)))

z_stat
p_value
\end{lstlisting}

\discussbox{The confidence interval and hypothesis test use the same sample, but they answer different questions. What question does each one answer?}

\writebox{At $\alpha=0.05$, do you reject $H_0$? Write your conclusion in context.}

\section*{Question 3. Simulate confidence interval coverage}

Create a hypothetical population where the true proportion is known. Then repeatedly sample from it and count how many confidence intervals contain the true proportion.

\begin{lstlisting}[language=R]
set.seed(213)

N <- 675000
pop <- c(rep("Cellphone", N * 0.9), rep("No Cellphone", N * 0.1))
true_prop <- mean(pop == "Cellphone")

true_prop
\end{lstlisting}

Complete the simulation below.

\begin{lstlisting}[language=R]
alpha <- 0.05
n <- 200
K <- 1000

num_of_simulation <- rep(0, K)
prop_est <- rep(0, K)
lower <- rep(0, K)
upper <- rep(0, K)
included <- rep(FALSE, K)

for (k in 1:K) {
  sampled_data <- sample(pop, size = ____)
  p_hat <- mean(sampled_data == "Cellphone")
  SE <- sqrt(p_hat * (1 - p_hat) / n)
  CI <- p_hat + c(-1, 1) * qnorm(1 - alpha / 2) * SE

  num_of_simulation[k] <- k
  prop_est[k] <- p_hat
  lower[k] <- CI[1]
  upper[k] <- CI[2]
  included[k] <- CI[1] <= true_prop & true_prop <= CI[2]
}

df_table <- data.frame(
  num_of_simulation = num_of_simulation,
  prop_est = prop_est,
  lower = lower,
  upper = upper,
  included = included
)

head(df_table)
mean(df_table$included)
\end{lstlisting}

\predictbox{Before running the simulation, about what proportion of 95\% confidence intervals do you expect to contain the true proportion?}

\writebox{What proportion of intervals contained the true proportion? Was it close to 95\%?}

\section*{Question 4. What changes when confidence level changes?}

Modify your simulation by changing \texttt{alpha} to 0.01 for a 99\% confidence interval.

\begin{lstlisting}[language=R]
# Change alpha to 0.01 and rerun the simulation.
# How does the coverage rate change?
# How does the interval width change?
\end{lstlisting}

\begin{center}
\begin{tabular}{|p{0.25\textwidth}|p{0.28\textwidth}|p{0.37\textwidth}|}
\hline
\textbf{Confidence level} & \textbf{Coverage rate} & \textbf{Typical interval width}\\
\hline
95\% & & \\[20pt]
\hline
99\% & & \\[20pt]
\hline
\end{tabular}
\end{center}

\debatebox{A higher confidence level sounds better. What is the tradeoff?}

\section*{Question 5. Build a two-way table}

For the chi-square test, use R's built-in \texttt{Titanic} table. The code below converts it to a data frame and summarizes survival by passenger class.

\begin{lstlisting}[language=R]
titanic_df <- as.data.frame(Titanic)

class_survival <- titanic_df %>%
  group_by(Class, Survived) %>%
  summarize(count = sum(Freq), .groups = "drop")

class_survival

tbl <- xtabs(count ~ Survived + Class, data = class_survival)
tbl
addmargins(tbl)
\end{lstlisting}

\discussbox{What would independence mean in this setting? In plain language, what would it mean if passenger class and survival were independent?}

\section*{Question 6. Calculate expected counts}

For each cell in a two-way table,
Expected count = (row total)(column total) / table total.

\begin{lstlisting}[language=R]
row_totals <- rowSums(tbl)
col_totals <- colSums(tbl)
total <- sum(tbl)

expected_tbl <- outer(row_totals, col_totals) / total
expected_tbl

all(expected_tbl >= 5)
\end{lstlisting}

\writebox{Do the expected counts meet the chi-square condition? How do you know?}

\section*{Question 7. Compute the chi-square statistic and p-value}

Compute the chi-square statistic manually, then use the chi-square distribution to find the p-value.

\[
X^2 = \sum \frac{(O-E)^2}{E}
\]

\begin{lstlisting}[language=R]
chi_stat <- sum((tbl - expected_tbl)^2 / expected_tbl)
df <- (nrow(tbl) - 1) * (ncol(tbl) - 1)
p_value <- 1 - pchisq(chi_stat, df = df)

chi_stat
df
p_value
\end{lstlisting}

\claimbox{At $\alpha=0.05$, I \underline{\hspace{2cm}} reject $H_0$. My conclusion in context is \underline{\hspace{7cm}}.}

\section*{Question 8. Try another pair of variables}

Repeat the chi-square workflow for \texttt{Sex} and \texttt{Survived}.

\begin{lstlisting}[language=R]
sex_survival <- titanic_df %>%
  group_by(Sex, Survived) %>%
  summarize(count = sum(Freq), .groups = "drop")

tbl2 <- xtabs(count ~ Survived + Sex, data = sex_survival)
tbl2

# Expected counts.

# Chi-square statistic.

# Degrees of freedom.

# P-value.
\end{lstlisting}

\writebox{Does survival appear independent of sex in the Titanic data? Support your answer with a p-value and context.}

\section*{Optional Challenge}

Write a reusable function that takes a two-way table and returns expected counts, chi-square statistic, degrees of freedom, and p-value.

\begin{lstlisting}[language=R]
manual_chisq <- function(tbl) {
  # Your code here.
}
\end{lstlisting}

\end{document}
